\documentclass[11pt]{article}
\usepackage[a4paper,margin=19mm,headheight=14pt]{geometry}
\usepackage{fontspec,xeCJK}
\IfFontExistsTF{Palatino}{\setmainfont{Palatino}}{\setmainfont{Latin Modern Roman}}
\IfFontExistsTF{Heiti SC}{\setCJKmainfont{Heiti SC}}{\setCJKmainfont{Noto Serif CJK SC}}
\xeCJKsetup{CJKmath=true}
\usepackage{amsmath,amssymb,enumitem,microtype,fancyhdr}
\usepackage[hidelinks]{hyperref}
\linespread{1.04}\setlength{\parindent}{0pt}\setlength{\parskip}{5pt}
\pagestyle{fancy}\fancyhf{}\fancyhead[L]{\small CCSGD / Reading guide}\fancyfoot[C]{\thepage}\renewcommand{\headrulewidth}{0pt}
\begin{document}
{\LARGE\bfseries CCSGD：中文阅读导读}\par{\small 与英文版逐节对应｜2026 年 10 月 11 日}\par

\section*{1. 论文解决什么问题？}
目标是优化干净输入下的目标函数 $F(\theta)=\mathbb E\ell(\theta;Z)$，实际只能观测 $W=Z+U$。即使输入噪声均值为零，$g(\theta;W)$ 通常仍有偏；降低学习率不能消除这个总体偏差。论文构造可直接执行、面向干净目标的随机梯度更新。

主体理论假设观测误差为已知半正定协方差的高斯噪声、潜在输入亚高斯，且每步样本新鲜并独立于历史。全纯延拓和输入增长条件针对梯度；凸性针对目标函数关于参数的几何结构。这两类要求需分别检查。


\section*{2. 如何分类问题？}
先分清修正问题与优化问题。高斯精确修正需要梯度关于输入的 entire 延拓及可积性，不需要凸性。要获得优化保证，先推导相对方差界，再按目标函数类别选用结果。

凸目标：定理 4.1 控制按步长加权平均的目标函数差。强凸目标：定理 4.2 证明到 $\theta^*$ 的几乎处处及均方收敛；定理 4.3--4.4 在额外局部条件下给出最后迭代和均匀平均的 CLT。非凸目标：定理 4.5 控制加权平均梯度平方，不保证到全局最优解的距离。

再按噪声分类。高斯噪声对应主体定理；均匀和对称双点噪声有明确的共轭扰动，见第 7.1 节，但共轭分布的指数尾部限制梯度增长。增长过快的 entire 梯度，以及任意非 entire 梯度，均不能自动获得保证。


\section*{3. 核心机制是什么？}
独立生成 $V_t\sim N(0,\Omega)$，令 $G_t=\operatorname{Re}g^{\mathbb C}(\theta_t;W_t+iV_t)$，再更新 $\theta_{t+1}=\theta_t-\eta_tG_t$。经典复数 corrected-score 恒等式使每步梯度精确无偏。本篇通过明确的矩估计将它接入优化理论，计算时无需先求闭式积分。

关键增长界为 $\|D_zg^{\mathbb C}(\theta;q)\|\le M(\theta)e^{a\|q\|}$，并要求 $M(\theta)^2\le A+B\|\nabla F(\theta)\|^2$。条件正交性分离原始采样噪声与修正误差，得到 $\mathbb E_t\|\xi_t\|^2\le\sigma^2+\rho\|\nabla F(\theta_t)\|^2$。高阶矩由增长界和分布尾部推出，不能仅从二阶矩假设推出。


\clearpage

\section*{4. 哪些结论便于复用？}
取 $\eta_t=c/(t+\gamma)^\alpha$，$1/2<\alpha\le1$。在规定的步长上界下，凸目标差和非凸加权平稳性指标在 $\alpha<1$ 时为 $O(T^{\alpha-1})$，在 $\alpha=1$ 时为 $O(1/\log T)$。强凸 MSE 在 $\alpha<1$ 时为 $O(t^{-\alpha})$；在 $\alpha=1$ 且 $2\mu c>1$ 时为 $O(t^{-1})$，临界和低于临界时分别有对数和较慢速率。

CLT 先使用收敛结论，再局部化，并用推导得到的 $(2+\delta)$ 阶矩验证 Lindeberg 条件；仍需协方差连续性和局部 Hessian 正则性。端点 $\alpha=1$ 要求 $c\lambda_{\min}(H)>1/2$，且 PR 协方差不同，不能直接代入区间内部的公式 $H^{-1}S_*H^{-1}$。

非高斯共轭情形中，$a<\lambda$ 足以保证修正恒等式，$2a<\lambda$ 足以完成矩控制，其中均匀噪声的 $\lambda=\pi/b$，双点噪声的 $\lambda=\pi/(2b)$。后一严格不等式还能容许适当的 $\delta>0$。多项式梯度容易满足输入矩条件，但仍需验证幅度控制及目标函数假设。


\section*{5. 实验说明什么？}
非线性特征回归关于系数是强凸的；附录 A.9 单独验证联合响应模型的矩条件。参数 MSE 拟合斜率为 $-0.978$，与 $O(t^{-1})$ 一致。虚线是指定更新规则的总体渐近平方偏差，不是每个有限时刻的下界。

正弦网络实验衡量干净输入的预测 MSE，属于经验研究。固定潜在训练样本每次使用时重抽观测噪声，无偏性针对干净的经验目标；不宣称网络全局收敛或预测误差趋零。到 200,000 步时，SIMEX 报告的 MSE 略小。WH 比较对象是其确定性式 (5)，不代表该文所有方法。


\section*{6. 建议阅读路径}
了解算法：第 3 节及算法 1，再读第 5 节例子。复用证明：引理 3.4 $\to$ 附录 A.2 $\to$ 第 4 节对应定理。理解 CLT：定理 4.2 $\to$ 附录 A.5 $\to$ A.6 或 A.7。实现非高斯修正：第 7.1 节及算法 2。核对实验范围：第 6 节、附录 A.9 和附录 B。

Agent 应先读五页 AI 版本，再用定理图中的稳定 ID 定位结果。原论文是数学叙述的权威版本；共轭变量判据属于外部依赖，作者提供的 6 Astra 评审未发现 main(6).pdf 存在大的数学缺口，但该文件与仓库快照的文件级对应尚未确认。

\end{document}
